%----------------------------------------------------------------------------------------
% Tailored CV — Research Engineer GenAI (Post-Training & Agentic AI)
% Fraunhofer IAIS, NetMedia, Dresden. TVöD E13, 2 years. Online portal.
% Lead: Agentic AI / multi-agent systems + foundation-model post-training + PyTorch/ML.
%----------------------------------------------------------------------------------------

\documentclass[10pt,a4paper,sans]{moderncv}

\usepackage[utf8]{inputenc}
\usepackage[T1]{fontenc}

\moderncvstyle{classic}
\moderncvcolor{blue}

\usepackage[scale=0.78]{geometry}
\usepackage{ragged2e}

%----------------------------------------------------------------------------------------

\firstname{Kundai Farai}
\familyname{Sachikonye}

\address{Auf der Lichtung 19}{82178 Puchheim, Germany}
\mobile{(+49) 1626386344}
\email{kundai.sachikonye@bitspark.com}
\homepage{github.com/fullscreen-triangle}{github.com/fullscreen-triangle}

\quote{GenAI research engineer: agentic and multi-agent systems, foundation-model
adaptation, and PyTorch ML --- with a builder's record of research turned into
deployed tools.}

%----------------------------------------------------------------------------------------

\begin{document}

\makecvtitle

%----------------------------------------------------------------------------------------
%	PROFILE
%----------------------------------------------------------------------------------------

\section{Profile}

\cvitem{}{
  Computational researcher and engineer working on \textbf{agentic / multi-agent
  AI}, \textbf{foundation-model adaptation}, and \textbf{ML in PyTorch}. I design
  conceptual frameworks and turn them into working, deployed systems. Strongly
  oriented to abstract and conceptual work, with a track record of research that
  ships.
}

%----------------------------------------------------------------------------------------
%	SELECTED GENAI SYSTEMS
%----------------------------------------------------------------------------------------

\section{Selected GenAI / Agentic Systems}

\cvitem{kwasa-kwasa / Turbulance}{
  \textbf{A declarative agentic-orchestration language.} Designed and implemented
  a programming language for multi-agent, evidence-bearing computation: an
  interpreter with confidence/argumentation semantics, multi-substrate delegation
  (LLM oracles, tools, Python specialists), and grounding of generated values.
  Rust core + TypeScript runtime; packaged as VSCode tooling incl.\ a language
  server. \href{https://github.com/fullscreen-triangle/kwasa-kwasa}{github.com/fullscreen-triangle/kwasa-kwasa}
}

\cvitem{four-sided-triangle}{
  \textbf{Metacognitive multi-agent pipeline.} An 8-stage RAG / reasoning pipeline
  with resource-aware orchestration, throttle handling, and a process monitor;
  Rust core (15--50$\times$), FastAPI + Ray + Docker.
  \href{https://github.com/fullscreen-triangle/four-sided-triangle}{github.com/fullscreen-triangle/four-sided-triangle}
}

\cvitem{Foundation-model adaptation}{
  \textbf{Post-training / fine-tuning.} Low-rank (LoRA) adaptation of pretrained
  models for domain tasks; knowledge distillation; a research line on training
  methodology (curriculum design / inverse curricula).
}

%----------------------------------------------------------------------------------------
%	EXPERIENCE
%----------------------------------------------------------------------------------------

\section{Experience}

\cventry{2021--present}{Independent AI / ML Researcher \& Engineer}{Bitspark GmbH}{}{}{
  Full-time design and implementation of agentic / multi-agent systems,
  foundation-model adaptation, and ML pipelines. From concept to deployed,
  documented, benchmarked tools (validated against reference data). Rust,
  Python/PyTorch, distributed/GPU.
}

\cventry{2019--2021}{Computational Science --- ML \& HPC}{Technische Universität München}{Munich}{}{
  Deep-learning pipelines (PyTorch, CNNs) at HPC scale; distributed processing of
  $>$100\,GB datasets on SLURM clusters; statistical evaluation and benchmarking.
  Taught at Master's level.
}

\cventry{2017--2018}{Data Pipelines}{Universitätsklinikum Hamburg-Eppendorf}{Hamburg}{}{
  Automated, reproducible high-throughput data-processing pipelines.
}

%----------------------------------------------------------------------------------------
%	EDUCATION
%----------------------------------------------------------------------------------------

\section{Education}

\cventry{2018--2019}{MSc Computational Biology}{Universität Tübingen}{}{}{
  Machine learning, statistics, and large-scale data analysis.
}
\cventry{2014--2017}{MSc Pharmaceutical Biotechnology}{HAW Hamburg}{}{}{}
\cventry{2011--2014}{BSc Biochemistry and Cell Biology}{Jacobs University Bremen}{}{}{}
\cvitem{}{\textit{Additional:} doctoral research, computational lipidomics, TUM
  (2019--2021; not completed).}

%----------------------------------------------------------------------------------------
%	TECHNICAL SKILLS
%----------------------------------------------------------------------------------------

\section{Technical Skills}

\cvitem{Agentic AI}{
  Multi-agent / agent orchestration systems; tool/LLM/oracle delegation;
  RAG and reasoning pipelines; evidence/grounding mechanisms
}

\cvitem{ML / foundation models}{
  PyTorch (CNNs, autoencoders, attention/transformers); foundation-model
  fine-tuning (LoRA); knowledge distillation; representation learning;
  evaluation \& benchmarking pipelines
}

\cvitem{Infrastructure}{
  GPU training; SLURM clusters; Ray, Dask; Docker, Kubernetes; FastAPI;
  cloud-portable workflows
}

\cvitem{Programming}{
  Python (primary) \& PyTorch, Rust, TypeScript, R, Bash, SQL;
  language/compiler design
}

%----------------------------------------------------------------------------------------
%	LANGUAGES
%----------------------------------------------------------------------------------------

\section{Languages}

\cvitemwithcomment{English}{Native / Fluent (C2)}{}
\cvitemwithcomment{German}{Conversational (improving)}{resident in Germany since 2011}

%----------------------------------------------------------------------------------------

\renewcommand{\listitemsymbol}{-~}

\end{document}
